\phantomsection
\addcontentsline{toc}{chapter}{LỜI CẢM ƠN}
\chapter*{LỜI CẢM ƠN}

Lời đầu tiên, nhóm nghiên cứu xin bày tỏ lòng biết ơn sâu sắc và chân thành nhất đến \textbf{PGS. TS. Nguyễn Thái Nghe}. Thầy đã tận tình hướng dẫn, định hướng khoa học, động viên và tạo mọi điều kiện thuận lợi nhất cho nhóm trong suốt quá trình triển khai và hoàn thiện đề tài \textit{“Hệ thống nhận diện thực phẩm và xác định giá trị dinh dưỡng ứng dụng Deep Learning”}. Những ý kiến đóng góp quý báu và sự chỉ dẫn nhiệt tâm của Thầy là nền tảng vững chắc giúp đề tài hoàn thành đúng tiến độ và mục tiêu đề ra.

Nhóm cũng xin trân trọng gửi lời cảm ơn đến Ban Giám hiệu \textbf{Trường Đại học Cần Thơ}, Ban Giám hiệu \textbf{Trường Công nghệ Thông tin và Truyền thông}, cùng \textbf{Phòng Khoa học, Công nghệ và Đổi mới sáng tạo} đã phê duyệt, hỗ trợ nguồn kinh phí nghiên cứu khoa học sinh viên và tạo môi trường học thuật lý tưởng để sinh viên phát huy năng lực nghiên cứu. 

Đồng thời, nhóm xin gửi lời tri ân đến toàn thể quý Thầy, Cô tại Trường Công nghệ Thông tin và Truyền thông – những người đã tận tâm giảng dạy, truyền đạt nền tảng kiến thức chuyên môn và kỹ năng thực hành vững chắc trong suốt những năm học vừa qua.

Cuối cùng, chúng em xin gửi lời cảm ơn chân thành đến gia đình, người thân và bạn bè đã luôn ở bên động viên, ủng hộ cả về vật chất lẫn tinh thần. Sự đồng lòng, nỗ lực và tinh thần làm việc nhóm nghiêm túc của tất cả các thành viên chính là động lực lớn nhất giúp đề tài đạt được kết quả như ngày hôm nay.

Mặc dù nhóm đã dành nhiều tâm huyết và nỗ lực để nghiên cứu, song do kiến thức và thời gian có hạn, báo cáo khó tránh khỏi những thiếu sót nhất định. Nhóm kính mong nhận được sự thông cảm, chỉ dẫn và những đóng góp quý báu từ Hội đồng nghiệm thu cùng quý Thầy Cô để đề tài ngày càng hoàn thiện hơn.

\vspace{1.0cm}
\begin{flushright}
    \textit{Cần Thơ, tháng 10 năm 2026}\par
    \vspace{0.2cm}
    \textbf{Nhóm sinh viên thực hiện đề tài}\par
    \vspace{0.15cm}
    \textit{(Ký và ghi rõ họ tên)}
\end{flushright}